% 原书第 577 页（PDF 第 590 页）
\BookPage{577}{590}
\BookHeading{乾隆年间的盐商与内帑}{乾隆年间的盐商与内帑}{赖惠敏*}
\BookSubheading{一、前言}

清代盐商的研究从日本学者佐伯富教授的论着《清代鹽政の研究》，何炳棣教授〈扬州盐商〉首开风气。近五十年，台湾学者徐泓教授更以毕生精力投注于盐业史研究，刘翠溶教授也曾研究这领域。大陆学者讨论盐业生产、行销、资本累积的论文更多，本文仅从〈内务府银库进项月折档〉〈内务府奏销档〉之记载，就乾隆年间盐商与内务府银库收入提出补充意见。

汪士信在〈乾隆时期徽商在两淮盐业经营中应得、实得利润与流向试析〉，估算两淮盐商支出的报效、借帑利息、余银和胶削，总数为138,179,100两，其中朝廷所得42,625,113两占18.68\%，盐政官员所得为95,482,800两，占41.84\%。以上四项中除了报效和借帑看起来和皇帝关系比较密切，余银和腹削似乎是官员所得。2不过，〈内务府银库进项月折档〉中却记载盐规银、引费银、铜觔水脚减半余平等，

\vfill
\begin{BookNotes}
\BookNote{*}{作者任职于中央研究院近代史研究所。}
\BookNote{1.}{Ho, Ping-ti (何炳棣) .1954. "The Salt Merchants of Yangchou: A Study of Commercial Capitalism in Eighteenth-century China."HarvardJournalof Asiatic Studies. 17(1/2): 130-168.巫仁恕译，〈扬州盐商:十八世纪中国商业资本研究〉《中国社会经济史研究》，1999年第2期，页59-76。佐伯富，《清代盐政研究》（京都大学东洋史研究会，1956)；徐泓，《清代两淮盐场的研究》（台北:嘉新水泥公司文化基金会，1972)；一，〈盐价、银钱比价：清代两淮盐商的成本、利润及其没落的原因〉，收入陈捷先、成崇德、李纪祥主编，《清史论集》（北京：人民出版社，2006），下册，页642-659; Ts'ui-jungLiu(刘翠溶), "Features of Imperfect Competition of the Ming-Ch'ing Salt Market,'The Second Conference onModern Economic History (I) (Taipei: The Institute of Economics,Academia Sinica,1989), pp.559-627.}
\BookNote{2.}{汪士信，〈乾隆时期徽商在两淮盐业经营中应得、实得利润与流向试析〉，《中国经济史研究》，1989年3期，页95-111；相关的论文还有：汪祟，〈明清两淮盐利个案两则〉《中国经济史研究》2000；3，页78-86、73；汪崇，〈乾隆朝徽商在淮盐经营中的获利估算〉《盐业史研究》，2000年第1期，页13-22；汪崇筼，〈清嘉道时期淮盐经营成本中的估算和讨论〉《盐业史研究》，2002年第1期，页7-17。}
\end{BookNotes}

% 原书第 578 页（PDF 第 591 页）
\BookPage{578}{591}
\noindent 都先后成为皇帝索要的目标。本文拟讨论的问题之一为「盐政官员所得」中，有哪些所得是盐政官员用来孝敬皇帝的。

其次，汪文中讨论乾隆下江南备公用款为927万两，为商捐部分。实际上，<内务府银库进项月折档〉记载似乎未达九百万。盐商支付内务府款项最多的是两淮盐引案发后赔补的提引余利，此案曾有日本学者讨论过，佐伯富在《清代鹽政の研究》书中称此案为乾隆三大案之一，若干官员收受盐商贿赂被处死。3滝野正二郎，〈清代乾隆年間における官僚と鹽商（一）、（二）一兩淮鹽引案を中心として〉探讨两淮盐政普福、高恒等官员涉及贪污被处分情形，他统计雍正至乾隆朝两淮、长芦、河东盐政的任期，从一年一任改成数年一任，内务府包衣久任「盐政」的理由是为了让他们熟悉与监督地方事务，同时维持稳定的盐税；这些职位掌握丰厚的利益，专差久任的结果，助长贪污舞弊的机会。4本文将利用档案探讨此案替皇室增添的收入。

中国第一历史档案馆藏〈内务府银库进项月折档〉（以下简称〈月折档〉）记载乾隆皇帝的银库，每年进项高达到一、两百多万两。5然〈月折档〉的记录并不完全，皇帝在奏折中朱批:「交圆明园」、「造办处」或「养心殿」时，这些银两就不进银库了。所以关于皇室从盐政所得的收入，还要参考一档馆发行的乾隆朝〈宫中殊批奏折〉「财政·盐务类」（以下简称〈朱批·盐务〉）、台北国立故宫博物院发行《宫中档乾隆朝奏折》，以及《中央研究院历史语言研究所现存清代内阁大库原藏明清档案》（以下简称《明清档案》），才能找到较完整的资料。6〈月折档〉始自乾隆十五年至六十年，其中缺四十二年，乾隆五十八年以后的资料也较

\vfill
\begin{BookNotes}
\BookNote{3.}{佐伯富，《清代鹽政の研究》（京都:京都大学东洋史研究会，1956)，页228-229。}
\BookNote{4.}{滝野正二郎，〈清代乾隆年間における官僚と鹽商一兩淮鹽引案を中心として（一)〉，《九州大学东洋史论集》第15期，1986，页83-106；〈清代乾隆年間における官僚と鹽商一兩淮鹽引案を中心として(二)〉，《九州大学东洋史论集》第22期，1994，页1-17。}
\BookNote{5.}{参见朱庆薇，〈内务府广储司六库月折档〉，《近代中国史研究通讯》，期34（2000），页201-}
\BookNote{206。}{}
\BookNote{6.}{乾隆朝〈宫中殊批奏折〉「财政·关税类」（北京：中国第一历史档案馆发行微卷，1987)；《中央研究院历史语言研究所现存清代内阁大库原藏明清档案》（台北：中央研究院历史语言研究所藏）（以下简称《明清档案》）；乾隆朝<内务府奏销档〉；〈内务府银库进项月折档〉（北京：中国第一历史档案馆藏)；国立故宫博物院编，《宫中档乾隆朝奏折》（台北:国立故宫博物院，1982)。}
\end{BookNotes}

% 原书第 579 页（PDF 第 592 页）
\BookPage{579}{592}
\noindent 少。此档案系由总管内务府大臣或银库郎中奏报盐政官员交银两数目，该款项可能和朱批奏折重复。〈朱批·盐务〉从乾隆元年至六十年，其中缺四十二至四十三、五十二至五十三年，台北国立故宫博物院的《宫中档乾隆朝奏折》刚好可以弭补。就档案的时间顺序来说，盐政官员奏折和银两送到京师，到银库点收经过一个多月后才在〈月折档〉出现。所以这三种档案必须核对才不致于重复计算。《明清档案》关于盐政的档案很少，若有也可能是重复的，真正能用的档案不到十件。

笔者将以上资料建档共一千五百笔左右，从中筛选出长芦、两淮、河东盐政交给皇帝的衙门经费、商捐、提引余利、利息等项，呈现许多有趣的结果。乾隆皇帝的盐政官员都是内务府包衣出身，皇帝说盐政官员养廉银太高，各级官员马上诚惶诚恐地裁减，将节省养廉银奉送皇帝；皇帝说宫廷日用增加，官员随即捧上十万两的「铜斤、引费、减半余平」、「外支银」、「盐规银」、「陋规银」等。况且，商人请借内帑银让他们「生意年年俱好，获利甚多」，好像达到扶植盐商的作用。等到盐商借欠皇帝的内帑超过一千万两时，盐商支付的利息比盐税还高，就摊派在众商身上，后来更成为嘉庆朝和道光朝的财政问题。

\BookSubheading{二、盐政衙门与皇室收入}

笔者曾探讨乾隆时期的税关与皇室财政，发现监督来自内务府包衣，他们为了孝敬皇帝，因而撙节税关衙门的开支，将银两解交皇帝，故有办公、节省、饭食、养廉银的名目。又，雍正朝实施耗羡归公之后，地方官员征收火耗的比例有三成、五成不等，乾隆皇帝下令改征一成，导致税关衙门用度不足，遂于正项下附征税银，故有平余、并平、罚倍、倾销、归公银等名目。在盐政衙门中也发现类似的情况，且有过之。以下分别叙述长芦、两淮、河东盐政衙门对皇室财政贡献。

\BookSubheading{（一）、长芦盐政}

户部山东清吏司掌盐课之政令，是盐政的最高行政机构。巡盐御史（或称盐课

\vfill
\begin{BookNotes}
\BookNote{7.}{参见拙作，〈清乾隆朝的税关与皇室财政〉，《中央研究院近代史研究所集刊》，期46 (2004年12月），页53-103。}
\end{BookNotes}

% 原书第 580 页（PDF 第 593 页）
\BookPage{580}{593}
\noindent 监察御史、盐政监察御史），乾隆后简称「盐政」，系户部差遣至各盐区的最高盐务专官，统管一区盐务。8长芦盐政管辖山东运司衙门和长芦运司衙门两个单位，分别解交内务府的银两。

1.山东运司衙门

山东运司衙门支援内务府的笔帖式银两，始自雍正年间。根据巡视长芦等处盐政伊拉齐奏：「山东盐务有额解笔帖式银，并盐政衙门应帮笔帖式银两二项，于雍正元年（1723）奉裁起，按年解赴内务府广储司库交纳在案。」山东运司一年内应解笔帖式银两三千六百一十七两一钱，并盐政衙门应帮笔帖式银两五千两，二项共银八千六百一十七两一钱，这制度一直维持到乾隆六十年。又有运司衙门的附属分司官员出缺或官员病故，必须扣存缺官养廉银两也解交内务府广储司库。9

其次，山东运司存有「盐当规例银」，作为赏捕弭盗之用。乾隆九年（1744）巡视长芦盐政伊拉齐奏称，该项经费自乾隆六年（1741）至九年除赏捕弭盗外，约存银3万两，加以本年解司银一万余两，约存银四万两，该地方公用既无所需，于存银内拨银二万两解交长芦盐政。此二万两应解交何处交收。朱批：「交海望，其水脚即在长芦支给。」10根据〈内务府银库进项月折档〉记载，山东盐当规银每年交给内务府约一万多两。

2.长芦盐政衙门

长芦盐税解到户部时有盘费、余平、饭银，以及镕铸元宝所需损耗，这些经费系由商人的盐税中附征，于每千两随解余平银二十五两、饭银七两。至雍正八年（1730）奉旨减去一半，每千两只随解减半余平银十二两五钱、饭银七两。雍正十二年（1734）复奉准部咨，凡有解部钱粮将减半余平并饭食，照例动支起解。至乾隆三年（1738）复奉上谕：从乾隆三年为始减半平余银两，一概停其解部，即存

\vfill
\begin{BookNotes}
\BookNote{8.}{陈锋，《清代盐政与盐税》（郑州:中州古籍出版社，1988)，页29-30。}
\BookNote{9.}{〈宫中朱批奏折〉，「财政类·盐务」，第0444-022号，乾隆四年十二月十日。}
\BookNote{10.}{〈宫中殊批奏折〉，「财政类·盐务」，第0449-015号，乾隆九年五月二十七日。}
\end{BookNotes}

% 原书第 581 页（PDF 第 594 页）
\BookPage{581}{594}
\noindent 贮本省司库。遇有荒歉，及裨益民生之务。11实际上，长芦盐政将此项银两奏报归公后，又说应否解内务府？皇帝乐见奴才的尽心，朱批：交内务府大臣等。如乾隆五年（1740）巡视长芦等处盐政伊拉齐奏：「长芦课银解部向收有余平、饭银及解官盘费等项，查尚有部饭银1,500两、解官盘费银250两、倾镕炉银460两。又商人例带添平银1,010两，计共节省存银3,220两，应一并详请奏报归公。奴才查俱收存运库，此项银两应否解赴广储司，抑解交何处谨请谕旨。」奉旨:「交海望」。12海望为内务府总管大臣，即该银入内务府广储司银库。

又，长芦盐政每年在春、秋两季前赴山东秤挚引盐，一年支给盘费三千两。官员若生病不能前往，改经费需扣存归公。再，山东查获私盐，原许兵丁自理私盐变价，还有扣存缺官养廉银两等，都由长芦盐政奏请解交内务府。13

长芦运使养廉银原来每年为10,000两，运同5,000两，运判2,000两。乾隆十二年（1747）裁减长芦盐务官员的养廉银，长芦运使酌留5,000两、运同2,000两、运判1,600两，其余银两悉行裁减归于公项。此项归公银和盈余，以及帮内务府笔帖式、护军校等都解交内务府，每年大约一万至二万余两。14

乾隆十三年（1748）暂署长芦盐政丽柱奏称，伊拉齐乾隆十二年七月十八日起至乾隆十三年三月二十一日离任，按日计算共应得养廉银一万三千二百九两四钱三分，除伊拉齐于上年十一月十二日两次取过银九千两、本年二月内批发天津理事同知修抹水师营城堡银三百两外，尚未及领取银三千九百九两四钱三分。又商人金义禀出，伊拉齐任内发商探办雀鸟银四百七十五两，应归入查封伊拉齐家产内，共

\vfill
\begin{BookNotes}
\BookNote{11.}{朱云锦，〈户部平余案略〉，收入贺长龄编，《皇朝经世文编》（台北：文海出版社，1979)，卷27，页993。}
\BookNote{12.}{〈宫中殊批奏折〉，「财政类·盐务」，第0444-029号，乾隆五年二月十八日。}
\BookNote{13.}{〈宫中殊批奏折〉，「财政类·盐务」，第0448-005号，乾隆八年闰四月二十九日。乾隆八年长芦盐政三保奏称，长芦盐政每年于春秋二季前赴山东秤挚引盐，蒙圣恩赏给盘费银三千两分做两次支领，一次应领盘费银一千五百两。本年春季奴才患病不能前去所有春季盘费银一千五百两理应归公，又盐政衙门书承人役等应得饭银，除照例发给外，尚扣存盘费银一百九十两四钱九分亦应归公。再查山东历年残票解费并给兵公费以及自理私盐变价等项存银六百九十五两九钱三分，以上三项共银二千三百八十六两四钱二分，俱非盐课正杂钱粮且并不应解部之项，均属闲款银两，解交内务府或径交何处收纳。朱批：「交内务府」。}
\BookNote{14.}{〈宫中朱批奏折〉，「财政类·盐务」，第0451-026号，乾隆十二年四月十三日。}
\end{BookNotes}

% 原书第 582 页（PDF 第 595 页）
\BookPage{582}{595}
\noindent 银四千三百八十四两四钱三分。乾隆十三年署理长芦盐政丽柱奏称，查津属各商向有每引交银二分五厘，以为公费按引合算每年共交银一万九千余两，名曰二分半。经前任盐臣三保于雍正十三年奏明「归公以来，凡有杂用取给于此，俱系各商自收自用，每年恣意开销殊多浮冒，凡不敷用复又滩派，众商并受其累。夫盐政盐道运同运判等关每年应得养廉银两足敷用度，凡修理衙署等项原应自行出资办理。」朱批：「俱交海望」。15

乾隆十四年（1749）奉上谕:「丽柱前因与商人往来交结，降旨询问伊具折覆奏等因着罚银一万两。官员罚俸银，自然也归内务府银库。」16该年，署理长芦盐政丽柱奏称，查办理盐务原系贸易生理与其派令官人经理不如招募民商为便，请将永庆号盐局21州县照吉安引窝之例，交与盐政招募殷实民商令其承办。今据运司叶昱详称募有商人沈朝安、查奕茂、卢敏功、朱立基、魏汝植俱系见在长芦办课，熟悉盐务殷实民商，尽可承办。并据商人刘公德、胡世亮、王克大、张锦具结公保倘有亏缺愿甘分赔。17

\BookSubheading{（二）、两淮盐政衙门}

两淮盐政衙门在雍正朝时候还没有任何经费解交内务府，到乾隆朝各种节省银、闲款、规费、外支银、铜斤水脚等名目，每年给内务府的银两不下于三、四十万两，这些经费都是从盐商输纳而来。两淮盐商除了正课之外，私费不赀，「遇一事即有一事之陋规，经一处即有一处之科派」，以下各项皆为盐商输纳内务府的各项银两。

1. 两淮盐政衙门节省银

两淮盐政衙门的节省银是来自商人按引缴纳，凡各项经费及京河各饷水脚饭

\vfill
\begin{BookNotes}
\BookNote{15.}{〈宫中殊批奏折〉，「财政类·盐务」，第0452-004号，乾隆十三年四月初四日。〈宫中朱批奏折〉，「财政类·盐务，第0452-008号，乾隆十三年四月初六日。乾隆十三年闰七月初九日海望奏请署长芦盐政丽柱解交银六千六百九十七两应交何处查收。奉旨着交造办处，钦此。}
\BookNote{16.}{〈宫中朱批奏折〉，「财政类·盐务」，第0453-017号，乾隆十四年十月二十日。丽柱解交银一万两应交何处查收。奉旨:「着交仓州，钦此」。}
\BookNote{17.}{〈宫中殊批奏折〉，「财政类·盐务」，第0453-035号，乾隆十五年二月十五日。}
\end{BookNotes}

% 原书第 583 页（PDF 第 596 页）
\BookPage{583}{596}
\noindent 食、官役养廉俸工，俱在此内支用。除动用之外其余存银两，向系盐政、运司各半分用，自雍正三年盐政噶尔泰始奏请节省充公，又裁减盐政运司两衙门书承各役工食银两，又裁减东司饭银等项，亦归于节省之内。在盐政经管者名为院节省；其在运司经管者名为司节省。雍正四年丙午纲遂入季报册报户部，自此历任盐政俱于考核时，将一纲节省银两具折奏明。18以上节省银照数解户部，雍正五年另给八旗养廉银24,000两。

乾隆三年两淮巡盐御史三保谨奏:「查两淮运库内存贮历年商捐各场火伏月费，及裁减盐务道书吏工食，并巡缉私盐兵役饭食，裁减盐义仓人役工食，以及淮南匣费等项共银十万九千一百八十五两八钱六分零。俱系存贮运库不入奏销，从未报部之项。奴才遵旨办理，当即提解到署，除将九千一百八十五两八钱六分解交苏州织造海保，请旨外，其余十万两奴才由水路解进，解交户部尚书海望收受。」19该奏折提及运库「不入奏销」的银两包括商捐的各场火伏月费、裁减盐务道书吏工食、巡缉私盐兵役饭食、裁减盐义仓人役工食、淮南匣费原贮存于运库，乾隆三年另有解交苏州织造局九千余两。

自乾隆五年以后，除两淮盐政衙门的节省银改拨四万两至内务府，据盐政三保的奏折云:「盐政衙门各项节省等银共十一万四千九百九十两二钱九分，内除例应支用水脚等银二万八千一百二十九两三钱八分，总理盐政督臣支养廉银三千两，奴才支用养廉银一万五千两，又帮贴总河臣养廉银七千两外，共节省银六万一千八百六十两九钱一分。谨将收支节省各数开折恭呈圣览。再据运使徐大枚详称该司衙门节省等银共七万一百五十两一钱五分零，内除例应支用薪俸等银一万九千四百二十二两六钱三分，实节省银五万七百二十七两五钱二分。遵例详请代奏等情前来理合一并开折带奏恭呈圣览。在查前项节省银两二共十一万二千五百八十八两四钱于本年接奉部文令于节省银内起解八旗养廉银二万四千两，实存节省银八万八千五百八十八两四钱合并陈明，伏乞皇上睿鉴施行谨奏。朱批：

\vfill
\begin{BookNotes}
\BookNote{18.}{〈宫中殊批奏折〉，「财政类·盐务」，第0443-039号，乾隆四年六月初四日。}
\BookNote{19.}{〈宫中殊批奏折〉，「财政类·盐务」，第0443-016号，乾隆三年三月十八日。}
\end{BookNotes}

% 原书第 584 页（PDF 第 597 页）
\BookPage{584}{597}
\noindent 「交海望四万两，其余交户部。」20乾隆十三年（1748）以后两淮盐政、盐运司的节省银全部解交内务府，每年都在十余万两。

乾隆十三年以后节省银项下增加倾销银，此因盐税起解户部时必须镕铸为元宝，若改拨其他地方之协饷，不须镕铸成元宝，该倾销银归入节省银，故乾隆中期后节省银多达十七、八万两。根据两淮盐政吉庆奏折，解部课银俱请元宝商人交易银两不能足色，一倾元宝必须加色、加平及火工饭食，当据运司详定每两需银0.033两，彼时两淮每纲额解正杂盐课2,678,000两，约计起解外省协饷1,07,.000余两，仍用小锭。起解京饷1,600,000两，及每年解部饭食并减半余平银约80,.000两，俱倾元宝，统计倾宝火耗公费等银约55,440两，令商带完应用。吉庆逐一核查卷案，自乾隆六年至十二年按每纲55,440两计，共应征银388,.080两，历年征用过349,.064.16两，尙短征银39,.015.84两，吉庆奏请自乾隆十三年（戊辰纲）为始，嗣后每纲照额征收倾费银55,440两存贮运库，视户部拨京饷多寡，照例支销倾费，若京饷拨不及数所有余存倾费汇入节省项下一并起解。21

\begin{BookNotes}
\BookNote{20.}{〈宫中朱批奏折〉，「财政类·盐务」，第0444-036号，乾隆五年三月二十七日。}
\BookNote{21.}{〈宫中殊批奏折〉，「财政类·盐务」，第0452-002号，乾隆十三年三月二十一日。}
\end{BookNotes}

\BookArt{p597-01.jpg}{0.990}{0.370}


% 原图说　资料来源：〈宫中朱批奏折〉，「财政类・盐务」，第0444-012号，乾隆四年九月二十六日

% 图表 1 的文字辨识草稿，行列按原图查看。
% 支出(两)
% 收入(两)
% 运司三分司各首领坝所各场员役俸
% 6,912.62
% 戊午纲带行戊申经费脚费各项节省
% 27,989.73
% 薪工食心红纸张等支给共银
% 等银
% 7,219.32
% 2,234.37
% 丁巳纲水脚等项共余银
% 四季扣解缺官俸工共银
% 1,180
% 淮南匣费节省官役等项共
% 奖赏急公商人花红联匾酒席共银
% 22,377.8
% 运使徐大枚支解节省银
% 2,779.63
% 捐解扬城土税等银
% 629.2
% 置备执事修理衙署考试券资节赏等
% 支放杂项并顺挑河并余平参吏公堂
% 2,099.17
% 3,301.28
% 等项节省
% 582.19
% 前任运使苏霖泓支解节省银
% 3,516.12
% 捐解塘兵工食
% 恩赏养廉银（运使事盐务道徐大枚
% 6,000
% 支放杂项并顺余平参吏公堂等项节
% 3,103.69
% 支用银3,430.77两前任运使苏霖泓
% 支用2,569.23两）
% 70,287.61
% 19,637.55
% 共银
% 资料来源:〈宫中殊批奏折〉，「财政类·盐务」，第0444-012号，乾隆四年九月二十六日

% 原书第 585 页（PDF 第 598 页）
\BookPage{585}{598}
2.盐政衙门闲款、屯费、规费

盐政衙门的闲款来源之一为因商捐河工剩余银两。乾隆四年两淮盐政武备院卿三保谨奏：「查戊午纲（乾隆三年）杂项内有众商带纳工项银二万两，此项银两虽完在库实系无着闲款。可否照去年起解闲款之例交与内大臣海望抑或解交何处？朱批:「解交海望为是」。22

第二种闲款来自补平银等。乾隆八年两淮盐政准泰奏：「两淮起解户部盐课向令商人按每千两缴补平银十五两为银库添平之用，有壬戌纲秋拨后续收钱粮内，先准部咨存拨淮徐等属赈务银六十万两，因非解部盐课所有补平银九千两，系属每年闲款外之闲款，理当归公，应否解交海望查收。」23

第三种闲款是累积各年积存的闲款。乾隆九年两淮盐政准泰谨奏：「两淮盐务所有运库每纲积存闲款银两均于奏销盘库后，查明实数奏解内大臣海望查收，自乾隆八年三月起至本年二月底积存闲款一万二千九百一两二分，应否仍解海望查收？」朱批：「交海望」。24

第四种事办贡剩余的银两。乾隆九年两淮盐政准泰谨奏：「两淮盐差理应备办贡物，奴才于乾隆五年八月到任查知两淮者，奴才伏思交商备办殊非谨敬之道，因抵任伊始循例在于闲款内动支银四千九百四十九两恭办。于乾隆六年六月二十七日附折奏明，嗣后闲款银两俱未动支，历年具奏请旨尽解交海望查收。但每年备办贡物约需银一万二千两，养廉银不敷支用，查从前商人备办原系商捐银两，是以奴才将闲款奏解之后一切备办价脚酌照一万二千两之数按行盐纲分陆续行文盐运司之取商捐银两核实恭办以省糜费。今蒙皇上天恩署理安徽巡抚核实辛酉纲（乾隆六年）至甲子纲（九年）除实在支用价脚外，尙存剩银六千两现贮运库应移交新任盐政吉庆听其自行接办。」25

屯费、规费的由来系自乾隆十五年，署理两淮盐政普福奏折:「两淮运盐有屯

\vfill
\begin{BookNotes}
\BookNote{22.}{〈宫中朱批奏折〉，「财政类·盐务」，第0443-027号，乾隆四年二月二十四日。}
\BookNote{23.}{宫中朱批奏折〉，「财政类·盐务」，第0448-043号，乾隆八年十二月十四日。}
\BookNote{24.}{宫中殊批奏折〉，「财政类·盐务」，第0449-011号，乾隆九年四月十二日。}
\BookNote{25.}{〈宫中殊批奏折〉，「财政类·盐务」，第0449-017号，乾隆九年七月初一日。}
\end{BookNotes}

% 原书第 586 页（PDF 第 599 页）
\BookPage{586}{599}
\noindent 船公费一项，向系留为巡缉之费，续因确查功效不许冒滥。至乾隆十五年，经吉庆查明积余银两，一并奏请充公请屯费、规费银解交内务府」吉庆奏称：「两淮运盐屯船昔有桥坝等微员规费，嗣经裁革，令商于船价内照数扣出充公名曰屯船充公，归入巡费，留为各隘口缉私之用。奴才莅任以来凡隘口必确核其有无实效，巡役必稽查其有无功绩，不任冒烂虚糜。旧设巡费已足敷用，是以前项屯船充公现积有银二万二千六百八十八两，又商人运盐出场经过阙坝等处验引扑戳向有些须规费，原系各该微员之陋规。奴才清查积弊饬令廉隅自爱各该员不敢私肥己橐，具已报出充为公用，亦有银一万三百五十两，以上屯费、规费二项共银三万三千三十八两，均不在正杂钱粮及节省款项之内，亦无须报部，可否解交海望以充公用。」26

由此可知，屯费是屯船公费，而规费是桥坝微员所征规费，这两项每年仅一二万两。

3.铜斤、引费、减半余平

铜斤水脚饭食属于两淮的杂课名目，又有两淮盐商致送上下两将巡抚司道，以及府厅州县各官规礼一项，以淮商名义致送称为「盐规银」，又以口岸卖商名义致送称「引费银」。减半余平的来源如长芦盐政衙门，盐税中每千两征减半余平银12.5两、饭银7两。这三项银缴交内务府的由来是在乾隆十三年时，吉庆奏称：「奉上谕两淮于闲款银内每年解银十万两交内务府备用。奴才已查明铜斤引费减半余平三款共成十万两，报部拨解。但铜斤一款虽经入册奏销，究与引费减半余平无异，原系盐务杂项并非正额盐课，可否仍遵前旨每年起解内务府以充公用。」27以上三项从乾隆十五年开始正式每年缴交内务府，银两为100,000两。

铜斤、引费、减半余平三项收入有时没交内务府，改发给淮商生息，每年的利息为12\%。乾隆三十三年两淮盐政尤拔世谨奏：「今查本年有戊子（三十三年）纲必带戊寅（二十三年）一分铜觔引费等银共十一万两，亦系例解内务府之项。奴才拟请循照前例即于此项内拨出银十万两抵做应领当本，给谙习营运商人领借按月一

\vfill
\begin{BookNotes}
\BookNote{26.}{〈宫中殊批奏折〉，「财政类·盐务」，第0453-037号，乾隆十五年二月二十二日。}
\BookNote{27.}{〈宫中殊批奏折〉，「财政类·盐务」，第0452-013号，乾隆十三年六月初七日。}
\end{BookNotes}

% 原书第 587 页（PDF 第 600 页）
\BookPage{587}{600}
\noindent 分起息。至铜斤引费项下尚余银一万两将来仍解交内务府。」28

4.吉安窝利、江西盐规引费

乾隆九年两淮盐政准泰谨奏：「吉安一府引盐原系侍郎永寿行盐地方，于雍正四年蒙特恩赏给怡贤亲王置办用度，嗣于雍正八年钦奉谕旨，交淮商黄光德办理。吉安府计额引五万一千五百四引，每引纳窝利银一两，共应纳窝利银五万一千五百四两，每年春秋二季交那运库起解王府。」29

乾隆十四年，因怡亲王府的家人朱隆阿预先向商人江启源等收取银两，乾隆皇帝愤怒之余，取消给王府的赏银，遂使该银归入皇帝荷包。两淮盐政高晋谨奏：「奉旨怡亲王家吉安地方行盐之成本内王子家现在既经亦有收过成本之处，将每年应得利银五万八千四百七十七两四钱内，于十三年十月内奉吉前院牌奉怡亲王寄字饬将己巳纲（十四年）应得窝利银二万八千两，交在扬家人朱隆阿赶办事务等因。当经传知承办商人江启源等陆续完缴过银一万八千两，交明朱隆阿收领，取具领状报明在案。其余应交银四万四百七十七两四钱应解交何处。」朱殊批：「仍交三河请旨」。30又奉旨:「每年所得利银内自当先尽应行解交之三万两收足，然后再将王府二万八千余两按季分之方为允当。今己巳纲应交窝利先经怡亲王寄字预支银一万八千两乃吉庆行文运使舒隆安催收给发，殊属不和，兹奉旨将每年应得利银，嗣后不必赏给与怡亲王，入于交官银内。原任两淮盐政吉庆名下照数赔补归款，以为滥行支发者戒。」31乾隆三十年（1765）两淮盐政高恒谨奏:「吉安一府窝产，江都、甘泉、仪征三县田房什物各项可得价二万七千七百一十两，照此贮入运库饬商一分生息每年可得息银三千三百二十五两。32自此之后，吉安窝利银成为皇帝的收入之一，每年约为五万余两。

乾隆二十八年（1763）以后，江西地区增加了「盐规引费」。两淮盐政高恒

\vfill
\begin{BookNotes}
\BookNote{28.}{〈宫中殊批奏折〉，「财政类·盐务」，第0463-006号，乾隆三十三年七月初四日。}
\BookNote{29.}{〈宫中殊批奏折〉，「财政类·盐务」，第0449-012号，乾隆九年五月初九日。}
\BookNote{30.}{〈宫中朱批奏折〉，「财政类·盐务」，第0452-050号，乾隆十四年六月初八日。}
\BookNote{31.}{〈宫中殊批奏折〉，「财政类·盐务，第0453-001号，乾隆十四年六月初八日。}
\BookNote{32.}{〈宫中朱批奏折〉，「财政类·盐务」，第0462-056号，乾隆三十二年五月二十四日。}
\end{BookNotes}

% 原书第 588 页（PDF 第 601 页）
\BookPage{588}{601}
\noindent 奏称:「二十八年等纲已经报运及未经报运各残引补完盐规引费，此项收银三万七百九十五两九钱乾隆三十年，江西盐规引费并吉饶引费除解交西盐道库银五千两外，存银三万四千四百四十三两五钱。又江西盐道解交历年积存盐规引费八万七千七百六十四两一钱。」33还有江西改运湖广的盐引，完纳口岸匣费、布税、充公等项余存银。乾隆三十一年两淮盐政普福奏称：「庚辰（二十五）、辛巳（二十六）、壬午（二十七）、癸未（二十八）等纲有派运江西改运湖广之残引，各商再运库一例完纳口岸匣费布税充公等项余存银四千八百八十六两，系多运楚盐不在口岸额费数内节省之项，奏解充公应否解交阿里衮。」34此项盐规引费等每年交给内务府银两约五、六万两，连吉安窝利银在内，江西湖广地区，每年给内务府的银两在十万两以上。

5.外支银、外支不敷银等

乾隆十一年（1746）两淮盐政吉庆谨奏：「两淮每纲设银外支银四万八千两，为奏考书工纸张新纲开桥开锁等事公费、赈恤火灾民窑、修理道路桥梁、给发沿江一带救生红船，并普济、育婴工食。庵观寺院香灯及助恤应酬等项支用，名曰外支银。向系按纲征贮运库，亦非各商缴送，内有官为给发者，运库给领。」3乾隆二十六年两淮盐政高恒谨奏：「军机大臣会同奴才议奏盐务章程内办理常贡每年仍以一万二千两作数如有交办理之事，酌定银数行司于外支银内动给等因，查每年外支银四万八千两，除办贡一万二千两外，实存三万六千两为，为各项公务支用，如有特交之事恐不敷用，据商人等秉称现因办理差物需领有提引可以通融，无须另议。应请自明年壬午新纲为始，每年以外支不敷银四万两备用。」36

外支银原本是地方的赈济公费，每年外支银四万八千两，乾隆二十六年从中拨出办贡银一万二千两，却因皇帝「特交之事」不敷四万两，故在外支银项目之外，尚有「外支不敷银」。在乾隆三十六年以后不论外支银或外支不敷银通通解交内务

\vfill
\begin{BookNotes}
\BookNote{33.}{国立故宫博物院编，《宫中档乾隆朝奏折》，第24辑，页589-590。}
\BookNote{34.}{〈宫中朱批奏折〉，「财政类·盐务」，第0462-015号，乾隆三十一年三月二十七日。}
\BookNote{35.}{〈宫中朱批奏折〉，「财政类·盐务」，第0451-017号，乾隆十二年二月初七日。「署理两淮盐政吉庆奏陈两淮盐规项下开销四款银两折。}
\BookNote{36.}{〈宫中朱批奏折〉，「财政类·盐务」，第0460-002号，乾隆二十六年六月十二日。}
\end{BookNotes}

% 原书第 589 页（PDF 第 602 页）
\BookPage{589}{602}
\noindent 府。此外，还附加一项叫「革除陋规银」的。该年两淮盐政李质颖奏称，将外支银四万八千两，又革除陋规银一万七千四百三十七两九钱，一并拨入外支项下共银六万五千四百七十三两九钱，存为办贡及特交器物装修等项之用。37王思治、金成基在〈清代前期两淮盐商的盛衰》提到两淮盐商每年办理玉器银四万两，常贡银六万五千四百余两」，38在宫中档殊批奏折这些银两的名称分别是外支不敷银、外支银和革除陋规银，一直到乾隆六十年这三项银两都还交内务府。

6.其他

其他的项目多半是不定额的，譬如盐政将河工剩余银两交内务府。乾隆十九年（1754）两淮盐政吉庆奏:「查雍正五年（1727）淮商公捐挑濬河道银两除动用外余剩银五千五百一十四两，又缴回核减等银七十六两，二共银五千五百九十两。前项余剩银两虽系留备各场挑河之用，但盐务些小工程商人随时捐办，而工费大者又动运库钱粮，是以二十余年并未用及此项存银。又乾隆二年兴挑盐运等河淮商，亦有公捐银两，除动支外，余剩银一万一千七百一十二两。又仪征地方系掣挚总汇之所，盐包狼籍，每有穷民扫积零盐，恐其私售，故设立官店饬商收买，转卖民食所得余利，除存店作本及杂用外，余银解交运库存贮，自乾隆十年至今陆续零星积存银四千五百八十六两，以上共银二万一千八百九十两，请解赴内务府海望交收」39

还有另一项是余银。王思治、金成基在〈清代前期两淮盐商的盛衰〉一文中讨论皇室、内务府与盐商之间重要经济关系，就是「余银缴公」问题。所谓的「余银」即「正额钱粮之外」的多余之银，这是一笔为数甚大的巨款。康熙年间，每年五、六十万两余银除皇帝用于织造、铜斤之外，其余由盐政上交内库，供皇帝奢侈挥霍。40不过，在乾隆朝余银解交内务府的次数好像很少，如乾隆三十二

\vfill
\begin{BookNotes}
\BookNote{37.}{〈宫中朱批奏折〉，「财政类·盐务」，第0451-017号，乾隆十二年二月初七日、第0466-010号，乾隆三十六年七月二十九日。乾隆30年两淮盐政高恒奏，自乾隆二十七年始每年将外支不敷银4万两存贮运库，凡有特交办理之事撙节动用。皇帝「特交办理之事」多半是成做玉器。}
\BookNote{38.}{王思治、金成基，〈清代前期两淮盐商的盛衰〉，《中国史研究》，1981年第2期，页66-84。}
\BookNote{39.}{〈宫中殊批奏折〉，「财政类·盐务」，第0455-025号，乾隆十九年七月十三日。}
\BookNote{40.}{王思治、金成基，〈清代前期两淮盐商的盛衰〉，《中国史研究》，1981年第2期，页66-84。}
\end{BookNotes}

% 原书第 590 页（PDF 第 603 页）
\BookPage{590}{603}
\noindent 年两淮盐政普福奏请余银等项银两解交内务府100,704.41两、三十九年解交余银38,744.89两。41到底余银从康熙朝到乾隆朝有什么样的变化，必须再找资料来探索。

以上乾隆朝六十年中，两淮盐政衙门给予内务府的收入共18,811,258.95两，大约占盐政官员所得95,482,800两之20\%。(参见图1）根据汪崇筼，〈乾隆朝两淮盐商输纳的检讨〉讨论两淮盐商的输纳之杂项、加斤余平、火足辛工三项，每纲共计银1,176,000两，「一切公私用度，此内已敷支给」。42每年一百多万两的地方经费中，约有310,000两需交给皇帝的，大约占26.66\%。

\BookArt{p603-01.jpg}{0.990}{0.332}

\BookSubheading{（三）、河东盐政}

1.河东盐政衙门节省经费

河东盐政衙门对皇室收入的贡献不大，其项目约与长芦相同。乾隆七年河东盐政尙琳奏：「河东一切盐务应用之处除额设充公银之外，另有附余钱粮亦充公用

\vfill
\begin{BookNotes}
\BookNote{41.}{〈宫中殊批奏折〉，「财政类·盐务，第0462-056号，乾隆三十二年五月二十四日；第0469-002号，乾隆三十九年十二月二十四日。}
\BookNote{42.}{汪崇，〈乾隆朝两淮盐商输纳的检讨〉，《盐业史研究》，2003年第2期，页27-35。陈锋引用盐务档案说乾隆33年的正课银为952,590两，杂课银已达1,1176,000余两，参见氏着，《清代盐政与盐税》（郑州:中州古籍出版社，1988)，页124。}
\end{BookNotes}

% 原图说　图1　两淮盐政衙门占总经费比例

% 原书第 591 页（PDF 第 604 页）
\BookPage{591}{604}
\noindent 一曰并余；一曰芦苇变价；一曰杂项节省。以上三项皆止详明盐臣批准动支，并未造册报部一并核销。并余银两向皆每平五十两加重二钱以为镕成元宝之费，雍正六年经前运使杨梦琰详明盐臣除充饷解费部院饭银渠堰工程、盐臣养廉仍行原封解发并无存留，其余支给只需库平所于每平加二钱之处悉行充公。自雍正六年至乾隆六年共收过银一千七百四十四两。自乾隆七年起仍合计每年所收公务官钱之处照平加增，一并造册送部候拨充饷。」43

芦苇变价缘盐池以内所产芦苇向皆在池官役私卖分肥，自雍正十三年经前盐臣孙嘉淦饬令变价贮库以充公用，因芦苇之丰歉视雨水之多寡，是以变价之数未经定额，乾隆六年共收银2,720两。共收钱33,700文，通数充公。至杂项节省银两或系公务内节省或系各兵役空缺一切杂项自应各归本款项下核销。尚琳称因系以公充公亦遂率行批准，先后共动支过银7,703两。乾隆七年起应令仍于本款之销，并册报部不得别附杂项名色。朱批：「此事不便交部，汝如是据实陈奏澄清办理固属可嘉，但恐太清则后难为继，汝其酌量办理或既称附余，至岁底奏明或交内务府或赏何人之处，请旨尚属可行之事耳，将此旨汝处记档亦不必咨部也，钦此」44皇帝既然如此暗示，吉庆便上奏折说：「奴才自行报解内务府查收」45

河东额引原417,440道，每引征收正杂盐课0.42两，每年额征课银171,720两，其余尚有征收公务银，为官役养廉银公事之用，因无定额亦无报销。雍正三年奏明额引除正课外每引额收公务银0.29两，每年额收公务银104,000两，以五万两充饷；其余五万四千两为官役养廉银及各项公事之用，是以每年正课充饷银共有220,000两。又因额引不敷民食，每年请领余引100,000道，每年约可收余引正课公务等银75,000两。雍正七年又因75额引不敷民食，每年请续增余引2、3万道至100,000道，每年约可收余引正课公务等银75,000两，亦系充饷之用。至乾隆五年续增余引至120,000道，又应多收余课等银15,000两、养廉归公等项又可收银15,000两，以上共销额余二项盐引共六十三、四万不等，银共四十余

\vfill
\begin{BookNotes}
\BookNote{43.}{〈宫中殊批奏折〉，「财政类·盐务」，第0447-002号，乾隆七年正月二十六日。}
\BookNote{44.}{〈宫中朱批奏折〉，「财政类·盐务」，第0447-021号，乾隆七年十一月二十九日。}
\BookNote{45.}{〈宫中殊批奏折〉，「财政类·盐务」，第0447-021号，乾隆七年十一月二十九日。}
\end{BookNotes}

% 原书第 592 页（PDF 第 605 页）
\BookPage{592}{605}
\noindent 万两。46

乾隆六年河东盐政尚琳奏，河南省南阳府属之唐县系食河东引盐，雍正十三年孙嘉淦奏令豫省商人公同协办每年缴余利银4,500两归公。凡土贩赴运领买每引一名仍召与商人交易之数纳银166两，除课额公务公费官钱于平之外，所于银两尽行归公收贮司库，每年于报销时具题并造册附款送部。47

2.唐县等余利银

乾隆七年河东盐政吉庆奏，据运使张任将委办商人范天锡办销乾隆七年分唐县额余二引并余利银两数目，具称办过额引40名107引，每名应得余利银54.7两，共该得余利银2,240.5两，又办过余引120名，每名应得余利银72.6两，共该得余利银8,717.1两，通计额余二引10,957两。查唐县余利银两向来报部归公原止银4,500两，其乾隆6年归公余利银除额缴4,500两外，另有3,600解交内务府。乾隆七年尚存余利银5,600较上年多银2,000两。48

乾隆八年河东盐政吉庆奏，看得襄城盐窝系张瑢、张公共之祖业，前年张瑢身故，妻孀弟幼一时乏人经营，课饷民食又刻不容缓，遂付托堂叔张珏将地方暂租与杨忠等供运虽议有压镇银6,000两，并每名长头银60两，杨忠等供称未实付压镇银，唯陆续付过长头银1,249.5两，应于张名下照追入官解交内务府。49

乾隆三十四年河东盐政固世衡奏，河东盐政应交商人节省银并裕州余利归公、唐县余利归公、赃罚余利归公，并借藩库随本减利等五项银两，拟由奏准之日起限五十日交到等因。查河东应交潞泽商人节省银20,000两、赃罚等银600至700两，并裕州余利归公1,000两，向于每年十月奏销后起解。唐县余利归公并带征余利归公，共银12,000两，向于奏销之次年三月内起解。各商原借藩库随本减利银约20,000两，向于一年期满后次年三月内起解。奴才与运使详查酌议后应请仍照历来办理成规，等五项银两，商人节省银并裕州余利归公、赃罚余利归公订于每年十

\vfill
\begin{BookNotes}
\BookNote{46.}{〈宫中朱批奏折〉，「财政类·盐务」，第0446-016号，乾隆六年三月二十一日。}
\BookNote{47.}{〈宫中朱批奏折〉，「财政类·盐务」，第0446-039号，乾隆六年十一月十六日。}
\BookNote{48.}{〈宫中殊批奏折〉，「财政类·盐务」，第0447-024号，乾隆八年二月二十九日。}
\BookNote{49.}{〈宫中殊批奏折〉，「财政类·盐务」，第0448-007号，乾隆八年五月初一日。}
\end{BookNotes}

% 原书第 593 页（PDF 第 606 页）
\BookPage{593}{606}
\noindent 月奏销后十一月起解，唐县余利归公、赃罚余利归公、各商原借藩库随本减利银于每年三月起解。今岁征收乾隆三十三年分商借藩库随本减利银27,000两，委员于三月解赴内务府。50

图2为乾隆年间长芦、两淮、河东盐政各项办公节省银等，每年解交内务府数量。

\BookArt{p606-01.jpg}{0.988}{0.351}

\noindent 三·盐商与皇室收入

\BookSubheading{（一）、盐商捐献}

乾隆九年淮南通河商人程可正等情愿公捐银16万两聊备内府公事支用。朱批:「传旨商人伊等急公献忱朕不容阻，但此等事以后可不必，但自循良完课即为良商矣。」奉旨:「雍和宫工程银五万、圆明园工程银五万，六万两交与仓州，钦此」51十二月淮南通河商人程可正等情愿公捐银15万两聊备内府公事支用。52

\vfill
\begin{BookNotes}
\BookNote{50.}{〈宫中朱批奏折〉，「财政类·盐务」，第0463-016号，乾隆三十四年四月初三日。}
\BookNote{51.}{〈宫中朱批奏折〉，「财政类·盐务」，第0449-007号，乾隆九年三月初四日。}
\BookNote{52.}{〈宫中朱批奏折〉，「财政类·盐务」，第0449-033号，乾隆九年十二月十二日。}
\end{BookNotes}

% 原图说　图2　盐政办公节省银

% 原书第 594 页（PDF 第 607 页）
\BookPage{594}{607}
\BookArt{p607-01.jpg}{0.990}{0.646}

大陆学者李克毅撰写〈清代盐商与帑银〉一文，认为导致盐商破产的原因并不是内务府发放的帑银造成的，而是皇帝与封建官僚的重重盘剥，盐商对皇帝的报效、索勒，使盐商们无力喘息。林永匡在〈乾隆帝与官吏对盐商额外盘剥剖析〉一文也认为乾隆帝与官吏对盐商，永无休止的需索再加上巨额的捐输、报效等，是造成两淮、常芦、河东等盐商衰落和破产的重要原因之一。53还有许多学者也有类似的看法，然就表1的数字看来，盐商对皇帝的报效不过15次，而且其中在乾隆三十

\vfill
\begin{BookNotes}
\BookNote{53.}{李克毅，〈清代盐商与帑银〉，《中国社会经济史研究》，1989年2期，页19-24；林永匡，〈乾隆帝与官吏对盐商额外盘剥剖析〉，《社会科学辑刊》，1984年3期，页90-95。}
\end{BookNotes}

% 原图说　表1　乾隆朝历年的商捐

% 图表 1 的文字辨识草稿，行列按原图查看。
% 表1乾隆朝历年的商捐
% 时间
% 商捐（两）
% 资料来源
% <殊批·盐务〉，编号0449-006、0449-033
% 乾隆九年(1744)
% 程可正等
% 310,000
% <朱批·盐务，编号0452-005
% 200,000
% 乾隆十三年(1748)
% 两淮盐商
% <朱批·盐务，编号0453-005、0453-024
% 乾隆十四年（1749）
% 30,000
% 程容德等
% 乾隆十六年（1751）
% 两淮盐商
% 400,000
% <朱批·盐务〉，编号0454-034
% 乾隆十七年(1752）
% <朱批·盐务〉，编号0454-047、0455-004
% 450,000
% 两淮盐商
% <朱批·盐务〉，编号0455-034
% 乾隆二十年(1755)
% 800,000
% 两淮盐商
% 乾隆二十二年（1757）
% 黄德源
% 1,000,000
% <殊批·盐务〉，编号0456-016
% 两淮盐商
% 乾隆二十五年（1760）
% 200,000
% 内务府银库进项月折档，25年2、12月
% <朱批·盐务〉，编号0460-047
% 100,000
% 乾隆二十七年（1762）
% 两淮盐商
% 《宫中档乾隆朝奏折》，第24辑，页445、
% 乾隆三十年(1765)
% 两淮盐商
% 460,090
% <内务府银库进项月折档〉，30年8月
% 乾隆四十五年（1780）
% <朱批·盐务〉，编号0471-029
% 河东盐商
% 100,000
% 乾隆四十九年(1784)
% 两淮盐商
% 1,000,000
% 内务府银库进项月折档〉，49年9月
% 乾隆五十一年(1786)
% <内务府银库进项月折档，51年12月
% 333,333
% 两淮盐商
% 乾隆五十三年（1788）
% 两淮盐商
% 333,333
% <内务府银库进项月折档〉，53年9月
% 乾隆五十六年（1791）
% 56,543.94
% 两淮盐商
% <内务府银库进项月折档，56年8、10月
% 56,543.94
% 乾隆五十七年(1792)
% 两淮盐商
% <内务府银库进项月折档〉，57年7、10月

% 原书第 595 页（PDF 第 608 页）
\BookPage{595}{608}
\noindent 年以前，这些负担尚不及提引余利或盐商长期借帑的利息。

商捐所占的比重

2,500,000

2,000,000

• ::: 1,500,000

0商捐总1,000,000

500,000

\begin{BookQuote}
0 e0 1740 1745 1750 1755 1760 1765 1770 1775 1780 1785 17901795 1800
\end{BookQuote}

西元

\BookSubheading{（二）、两淮盐引案与提引余利}

十八世纪中国人口迅速的增长，扩大食盐的消费市场，乾隆初便出现额引不敷民食的情况，自乾隆十一年两淮盐政吉庆建议预提盐引，以接济民食。两淮盐销售畅旺到前所未有的局面，但乾隆三十三年却发生两淮盐引案。日本学者佐伯富的看法是：「乾隆时代盐务极盛，额引畅销，不足者，预提翌年额引」。此案被称为乾隆三大案之一，若干官员收受盐商贿赂被处死。4根据〈内务府银库进项月折档〉的记载，此案不光是贪污案，而且商人赔补提引余利多达八百万两，通通都落入皇帝的荷包。

两淮盐引预提始自乾隆十一年，两淮盐政吉庆谨奏：

\begin{BookQuote}
奴才于乾隆十一年奏蒙恩准预提甲戌纲（十九年）盐引一百四十一万五千三百九十四引，以备畅销口岸，接济民食。嗣因每年提销所剩无几。又经奴才事先筹杂于乾隆十三年奏请照例预提备济，部议两淮盐引岁额不敷，令盐政先期约数奏请于次年额引内提出行销，引内注明预提字样，不足请增余剩缴部，其征课奏销悉照带乙提甲之例办理。查本年己巳纲（十四年）所领三十万
\end{BookQuote}

\vfill
\begin{BookNotes}
\BookNote{54.}{佐伯富，《清代鹽政の研究》（京都:京都大学东洋史研究会，1956)，页228-229。}
\end{BookNotes}

% 原书第 596 页（PDF 第 609 页）
\BookPage{596}{609}
\begin{BookQuote}
提甲引目业据各商捆配运销共纳正杂课银四十二万余两，商民两有利益，计自丙寅（十一年）、丁卯（十二年）、戊辰（十三年）、已巳（十四年）四纲已提销过一百二十万引提甲引目，仅剩二十一万五千三百九十四引，庚午（十五年）虽无闰月岁额之外尚须预为备济，所存提甲引目恐不敷销售，再将次年额引预提二十万道共成四十一万五千三百九十四引以备接济民食。55
\end{BookQuote}

由此奏折可看出，乾隆十一到十四年预提盐引120万引。余利银的来源则是因正引除了正杂课外，还加窝价银。在彰宝、尤拔世的奏折就说：「两淮纲引正引应完正杂课银二两六钱四分，此外尚有窝价一两。惟预提引张美引正杂钱粮不过一两五钱，不须另加窝价，并免纳杂费。（商人）预提引张系额外给予，领运其所得利息，亦系额外羡余。」56两淮盐引之正引纳正、杂、窝价银共3.6两，而预提盐引则缴交1.5两，所以每多2.1两的「余利」。这些提引余利银有部分做为地方公费，如普福奏称：「两淮每遇口岸畅销本纲额引不敷，俱循例奏请预提下纲引目给商领运，以资民食，商人亦沾余利并可为承办差务，以及地方公务之用。」地方公务，譬如：修茸淮北惠济祠、杨州府城隍庙各工程，江都、甘泉二县挑江修井添设防火水桶，并盐务寒冬收养穷民房舍等费等项。尤其像皇帝南巡沿途奖赏金牌银牌等都来自提引余利。57

但是，乾隆皇帝不管这项经费如何花用，令两淮盐商必须归还公款赔补10,200,800两。盐商提出赔补银两的困难所在：第一、历纲提引之总散各商事经二十余年，其中消乏存亡参差不一，统计历纲提引册内1,300余人，如今通河办

\vfill
\begin{BookNotes}
\BookNote{55.}{〈宫中朱批奏折〉，「财政类·盐务」，第0453-033号，乾隆十四年十二月二十六日。}
\BookNote{56.}{国立故宫博物院编，《宫中档乾隆朝奏折》，第31辑，页107。}
\BookNote{57.}{〈宫中朱批奏折〉，「财政类·盐务」，第0462-056号，乾隆三十二年五月二十四日。赏用金牌银牌如下：奖武金牌四十面重二十两，值银三百零四两奖武银牌一万五千六百面，重一万三千七百七十两养老银牌一万四千面，重四千二百二两六钱二两小锭四百五十锭，重九百两一两小锭五百二十六锭，重五百二十六两五钱小锭一千七百九十二锭，重八百九十六两三钱小锭二百七十锭，重八十一两二钱小锭五万三千三百八十五锭，重一万六百七十七两一钱小锭六百七十一锭，重七十两}
\end{BookNotes}

% 原书第 597 页（PDF 第 610 页）
\BookPage{597}{610}
\noindent 运现在现在只500余商。今700余人拖欠数百万两之公帑，现500余商，如仍按名勒追，必致大半无着。58第二、皇帝南巡的时候，盐商修葺园亭工程，如江广达、黄源德等修趣园、倚虹园、净香园、平山堂、观音山、九峰园等处，所费不眦。乾隆皇帝朱批却是：「一切工作悉皆取给公中，伊等因系经手之总商任意开销装饰，彼时既冒独立办差之名，滥膺恩赏，且以私家别业借口公事糜帑，增新仍得据为己有。及至事已败露所费之数，转令众商代为完缴。」59此指称盐商任意滥用公款办差，装修私家别业，事发后有转嫁众商之嫌。

林永匡根据《扬州行宫名胜全图》档案，发现扬州行宫合计共建造宫殿楼廊5,154间，亭台196座。楼廊注明为商人建造者3,981间，未注商人姓名者1,173间。亭台注明商人建造者160座，未注商人姓名者36座。这些行宫修缮完毕后，盐商们还得继续解囊购置宫中陈设景物，如高旻寺行宫、天宁寺行宫，商人们购置的古玩珍宝、花木竹石不可胜数；平山堂行宫中本无梅花，乾隆首次南巡时，盐商们耗银植梅万株，专供皇上观赏。60盐商出资建造皇帝行宫园林有注明标记，并不是皇帝所说:「以私家别业，借口公事糜帑。」

不论盐商过去将提引余利用在地方公费，或者如滝野正二郎所讨论官员从中取利，乾隆认定商人摊还历年提引银10,200,800两，从三十四年开始分期摊还。四十年，两淮盐政伊龄阿谨奏：淮南各商应缴历纲提引余利，共银10,200,862两，自乾隆三十三年十月起分限十年完缴。于戊子（三十三）、己丑（三十四）两限完过银2,270,548两。三十五年仰邀天恩展限六年续完过庚寅（三十五）一限银566,450两。三十六年内复蒙天恩再予展限三年，迨后辛卯（三十六）、壬辰（三十七）、癸巳（三十八）据遵加展年数匀分按限征完，每年分做两次起解。自三十三年起至三十八年共完解过银4,217,723两，尚欠5,983,138两。自甲午(三十九）7限起计13限，每限该交银460,241两。应解甲午首季银230,000两。61

巡视两淮盐政寅着奏称，商人共完过第7限4,677,965两，尙未完

\vfill
\begin{BookNotes}
\BookNote{58.}{国立故宫博物院编，《宫中档乾隆朝奏折》，第25辑，页731-733。}
\BookNote{59.}{〈宫中殊批奏折〉，「财政类·盐务」，第0463-017号，乾隆三十四年四月初七日。}
\BookNote{60.}{林永匡，〈乾隆帝与官吏对盐商额外盘剥剖析〉，《社会科学辑刊》，1984年3期，页91。}
\BookNote{61.}{〈宫中殊批奏折〉，「财政类·盐务」，第0469-015号，乾隆四十年五月十五日。}
\end{BookNotes}

% 原书第 598 页（PDF 第 611 页）
\BookPage{598}{611}
\noindent 银5,522,896两。又于四十年奉谕旨:「加恩展作二十限完纳，每限应完二十七万六千一百四十四两，商人应完四十一年丙申纲，自上年十月起至今年十月一年陆续完纳十四万两，未完纳一半银仍督商谨遵完纳。」2巡视两淮盐政伊龄阿奏称，乾隆四十七年奉上谕:「将未完银分二十一限缴完，今该商已完至第十三限，其余八限尙应缴银三百八十六万六千两。该商等按限完交已经过半，朕思恤商所以惠民，着加恩于未完银内豁免二百万两，其余一百八十六万六千两仍照原限完缴。」每年应交233,253.46两，分坐上下半年缴完。至乾隆四十九年奉旨:淮扬商众所欠1,632,774两，加恩全行豁免。63表2为两淮盐引案，商人历年缴交的提引余利银,共银6,510,491两。

\BookArt{p611-01.jpg}{0.990}{0.564}

\vfill
\begin{BookNotes}
\BookNote{62.}{国立故宫博物院编，《宫中档乾隆朝奏折》，第38辑，页717。}
\BookNote{63.}{清高宗敕撰，《清朝文献通考》（台北:台湾商务印书馆，1987)，卷29，页5123。}
\end{BookNotes}

% 原图说　表2　两淮盐商缴交历年盐引余利银

% 图表 1 的文字辨识草稿，行列按原图查看。
% 表2两淮盐商缴交历年盐引余利银
% 提引余利（两）
% 时间
% 资料来源
% 乾隆三十四年（1769）
% 1,270,548
% <内务府银库进项月折档，三十四年九、十一月
% 乾隆三十五年（1770）
% 10,00,000
% <内务府银库进项月折档〉，三十五年十月
% 乾隆三十六年（1771）
% 705,000
% <内务府银库进项月折档，三十六年六、七月
% 乾隆三十七年(1772)
% <内务府银库进项月折档〉，三十七年八月
% 230,000
% 乾隆三十八年(1773)
% 630,241.4
% <内务府银库进项月折档〉，三十八年闰三、七、十月
% 430,241.4
% 乾隆三十九年(1774)
% <内务府银库进项月折档〉，三十九年七月
% 乾隆四十年(1775)
% 430,241.4
% <内务府银库进项月折档>，四十年五、八月
% 乾隆四十一年（1776）
% 340,241.4
% <内务府银库进项月折档，四十一年五、九月
% 140,000
% 乾隆四十二年（1777）
% <内务府银库进项月折档>，四十二年五月
% 乾隆四十三年（1778）
% 246,144.84
% <内务府银库进项月折档，四十三年六月
% 246,144.84
% <内务府银库进项月折档，四十四年五、十一月
% 乾隆四十四年（1779）
% 106,144.84
% <内务府银库进项月折档>，四十五年六月
% 乾隆四十五年(1780)
% 乾隆四十六年(1781)
% 140,000
% <内务府银库进项月折档〉，四十六年十一月

% 原书第 599 页（PDF 第 612 页）
\BookPage{599}{612}
\BookArt{p612-01.jpg}{0.990}{0.134}

两淮盐引案商人应赔一千多万两，经皇帝加恩豁免，只交纳六百万余两。好像皇帝特别宽容，不过，在乾隆三十六年又开始实施预提盐引制度，在八年中补足皇帝豁免的部分。两淮盐政李质颖谨奏：「查提引余息惟视窝利之厚薄以定数目，乾隆三十三年尤拔世请题已丑纲引二十万道，每引缴银二两六钱共合五十二万两，今案出售情形，每引应缴余利二两八钱合计共缴银五十六万两。」4至乾隆四十年，每引改征余利3两，如此一来盐商每年缴给内务府的余利增为60万两。（参见表3、图4）根据陈锋的研究乾隆中期以后由于引额的壅积和报效数额增大，积欠亦日趋严重。如乾隆三十年的钱粮奏销册就出现「带征」戊寅（二十三年）、甲寅（二十九年）引课。此后，在奏销盐课时，「压征」、「接征」、「带征」等名目相继出现。两淮盐政李星沅指出:「带征远年旧欠，责成本愈重，商力不支。」65商力不支，却又增加提引余利造成商人的双重负担。

\BookArt{p612-02.jpg}{0.990}{0.373}

\vfill
\begin{BookNotes}
\BookNote{64.}{〈宫中殊批奏折〉，「财政类·盐务」，第0466-048号，乾隆三十七年三月二十九日。}
\BookNote{65.}{陈锋，《清代盐政与盐税》，页257-258。}
\end{BookNotes}

% 图表 1 的文字辨识草稿，行列按原图查看。
% 乾隆四十七年（1782）
% 106,144.84
% <内务府银库进项月折档〉，四十七年四月
% 〈内务府银库进项月折档〉，四十八年五、七、十月
% 乾隆四十八年（1783）
% 376,144.84
% 乾隆四十九年（1784）
% 113,253.46
% <内务府银库进项月折档〉，四十九年七月

% 原图说　表3　盐引余利银

% 图表 2 的文字辨识草稿，行列按原图查看。
% 表3盐引余利银
% 年代
% 资料来源
% 银两
% 乾隆三十七年(1772）
% 560,000
% <内务府银库进项月折档〉，三十七年八月
% <内务府银库进项月折档〉，三十九年七月
% 560,000
% 乾隆三十九年（1774）
% 乾隆四十年（1775）
% 600,000
% <内务府银库进项月折档〉，四十年八月
% 600,000
% 乾隆四十一年(1776)
% <内务府银库进项月折档〉，四十一年十月
% 600,000
% 乾隆四十三年(1777)
% <内务府银库进项月折档〉，四十三年十一月
% 600,000
% <内务府银库进项月折档〉，四十五年十一月
% 乾隆四十五年(1778)
% 600,000
% 乾隆五十八年(1779）
% <内务府银库进项月折档〉，五十八年五月
% 600,000
% <内务府银库进项月折档，五十九年八月
% 乾隆五十九年（1780）

% 原书第 600 页（PDF 第 613 页）
\BookPage{600}{613}
\BookArt{p613-01.jpg}{0.976}{0.429}

\BookSubheading{（三）、盐商借款}

盐商向内务府借内帑是康熙时代就开始的，乾隆朝盐商借内帑也十分常见。几年前笔者曾探讨商人王至德、范清济家族借内帑，当时我还没看到银库月折档，未能统计发商生息的数目。66从档案中可知，盐商向内务府借银两以长芦盐政占多数，其次是两淮盐商。（参见图5)

1. 长芦盐商借款

①盐政运库之公款

盐商向盐政衙门的运库借款为相沿的陋习，凡有商人缺乏起解钱粮即向运库暂为垫借陆续完交归款。67乾隆二十年署理长芦盐政官着奏称，乾隆十三年钦奉谕旨，令长芦于盈余闲款内每年解到五万两以备行幸赏赐之用。经前任盐政丽柱以芦

\vfill
\begin{BookNotes}
\BookNote{66.}{参见拙作，〈乾隆朝内务府的当铺与发商生息（1736-1795）〉，《中央研究院近代史研究所集刊》，28期（1997年12月），页133-175。}
\BookNote{67.}{〈宫中殊批奏折〉，「财政类·盐务」，第0461-008号，乾隆二十七年七月初四日。}
\end{BookNotes}

% 原图说　图4　两淮盐商缴交提引余利

% 原书第 601 页（PDF 第 614 页）
\BookPage{601}{614}
\BookArt{p614-01.jpg}{0.990}{0.373}

\noindent 东公费银32,500两堪充闲款，不敷银17,500两商人情愿捐足共凑50,000两。自乾隆十四年六月起至二十年五月止，营运六年可得本利银486,950两。68本利银内存留银417,000两，仍交商营运，其余69,950两解送内务府以为公用。嗣后每年应所得利银50,000两解交内务府。69长芦盐政衙门的经费可能是借给盐商，故其衙门经费缴交内务府的数量不及两淮盐政。

在官员每年的奏销册中，报告将该年地方经费用途，其中包括商人借款数量，分为已还或未完。乾隆二十七年长芦盐政达色奏称，乾隆二十六年分正杂课94,524.23两，平饭等银14,412.2两，共应108,936.43两，除盘查现存库银41,320.57两外，商借未完银72,473.86两，各商急公借用办理庆典皇差银两，除已完外尙有未完26,747.97两。各商急公之项现在陆续完外，各商垫借过银45,725.89两。内前任病故运使王图炯代商借过银27,616.93两、副运使印务运同马德生代商借过18,108.96两，二项借银内有正课银40,867.89两，闲款

\vfill
\begin{BookNotes}
\BookNote{68.}{〈宫中朱批奏折〉，「财政类·盐务」，第0455-015号，乾隆十八年九月二十九日。}
\BookNote{69.}{〈宫中殊批奏折〉，「财政类·盐务」，第0455-035号，乾隆二十年十月十七日。}
\end{BookNotes}

% 原图说　图5　盐商借内帑之利息

% 原书第 602 页（PDF 第 615 页）
\BookPage{602}{615}
\noindent 银4,858两。70达色的奏折提到商人借正课银，正课银不是应该缴交户部？又为何借户部的银两，利息又归皇帝所有？

②当铺

乾隆十二年以后内务府先后开设当铺共26座，估计成本银约726,548两。71其中有三十多万两交由长芦盐政发商生息，根据乾隆三十四年六月巡视长芦盐政高诚奏称，赏借当架本银352,612.46两，每年应交利银42,313.5两。72遇有闰月多1\%的利息。

另外，盐商也有个别向内务府借当架银。乾隆三十四年长芦盐政高诚奏称，商人王至德借过万成当官本银150,000两，又欠长芦运库官本银50,055两，共欠官本银200,055两，令该商每年本利兼交，利随本减，自丁亥年（三十二）起至壬寅年（四十七）止，分限十六年本利清还。王至德名下丁亥年应交利银24,000两，拔还本银10,055两，止交过利银18,000两。戊子年应交利银22,800，拔还本银10,000两，只交过利银18,000两，共应交利银46,000余两，今仅收过利银36,000两，较之原奏欠银10,000余两。73

不过，多数盐商借内帑通常不只当本一项，像王镗的借款就包括五项。乾隆二十六年长芦盐政金辉奏称，王镗引名王克大行销河南彰德府属之安阳、林县二处引地，每年共配运19,121引，每年应交课银10,326两。又领借当本生息银、赏赉生息银、内府帑利银、水师营生息、莲池书院生息银，以上五项共领借本银77,873.2两，应交利银9,344.7两。74王镗应交的正课银和他借帑银交的利息相差不远，可见长芦盐商负债累累。

\vfill
\begin{BookNotes}
\BookNote{70.}{〈宫中朱批奏折〉，「财政类·盐务」，第0461-009号，乾隆二十七年七月初四日。乾隆27年长芦盐政达色奏称，查商人借用长芦盐政库银两借办庆典银1,427.8两，借办皇差银53,585.5两，共借用银64,013.3两，今查两芦商人已完银37,265.33两，未完银26,747.97两。}
\BookNote{71.}{参见拙作，〈乾隆朝内务府的当铺与发商生息（1736-1795)〉，《中央研究院近代史研究所集刊》，28期（1997年12月），页133-175。}
\BookNote{72.}{〈宫中殊批奏折〉，「财政类·盐务」，第0463-042号，乾隆34年6月23日。}
\BookNote{73.}{〈宫中朱批奏折〉，「财政类·盐务」，第0463-048号，乾隆34年6月30日。}
\BookNote{74.}{〈宫中朱批奏折〉，「财政类·盐务」，第0460-005号，乾隆26年6月28日。}
\end{BookNotes}

% 原书第 603 页（PDF 第 616 页）
\BookPage{603}{616}
③代认革商欠银

长芦盐商因欠课被革退，由众商中选出财力雄厚的商人替他还帑息，「先将参商名下，所欠有利之帑，令其换具领状，认领输息，以抵引窝价值。其余旧欠课项，除参商家产盐价变抵外，不敷银两，核其旧欠银若于，计其所得利息多寡，酌定年限，并令新商代为完纳。」新商替旧商完纳帑息以永庆号最为著名。永庆号原系朱立基经营，因他亏空，由王得宜（起凤）接办永庆号盐务，永庆号的引地遍及直豫二省21州县，窝价银达46万两。又朱立基所积存盐包、房屋、家伙值银32万两，实在成本银819,672两。王起凤等承办每年以一分起息，应交库平色息银91,082两，定于乾隆十四年十二月二十八日自赴广储司交纳。76王起凤、王世荣、王珮三代都还经营永庆号，到乾隆晚期，永庆号缴交的利息达三百万两，皇帝称他「实属尚有天良，长久出力之商」。77

王珮在乾隆五十七年因雨泽愆期，销盐短少，请求内帑赏借80万两，情愿加倍缴银160万两，分限十六年交清。7乾隆五十八年，长芦盐政穆腾额奏称，王珮经营的永庆号官引地，及自置引地共50余处大半皆在河南，共窝价银650,000余两。各引地存积盐包，及存坨生盐共47万余包，共估值银1,660,000余两。除借过帑本扣抵外尚余银1,200,000两，所有该商京津当铺产业未经算入，通计抵借帑银800,000两不致有缺。可见盐商向皇帝借钱，得有雄厚产业作保来抵借帑银，不致有缺。 2

另一个商号为常茂号，原系王镗经营，他也曾欠永庆号银两，乾隆二十六年王

\vfill
\begin{BookNotes}
\BookNote{75.}{黄掌纶，《长芦盐法志》，收入《续修四库全书、史部、政书类840》（上海：上海古籍出版社据嘉庆十年刻本影印，1997），卷16，页23。}
\BookNote{76.}{〈宫中殊批奏折〉，「财政类·盐务」，第0453-015号，乾隆十四年十月十五日；《宫中档乾隆朝奏折》，第23辑，页309-311。长芦盐政高诚奏称，商人王得宜本名王起凤，前于内务府呈称伊自乾隆二十七年冬底承办永庆号盐务。缘前商朱立基亏缺成本银二十五万两余，未曾交付，所有永庆号运费资本悉系王得宜自行垫办。今有应交永庆号上年帑利银八万八千余两，并分赔朱立基亏缺银二万两，及另欠内务府库货价银十二万三千两，以上三项共银二十三万两。如按限交纳惟恐后力不继，现在备银二万五千两先行交纳，其余未完银二十万两，自三十年为使每年带交银二万五千两，分限八年带完。}
\BookNote{77.}{〈宫中殊批奏折〉，「财政类·盐务」，第0479-043号，乾隆五十八年十二月初三日。}
\BookNote{78.}{〈宫中殊批奏折〉，「财政类·盐务」，第0479-015号，乾隆五十八年五月二十八日。}
\end{BookNotes}

% 原书第 604 页（PDF 第 617 页）
\BookPage{604}{617}
\noindent 镗病故，伊子王德法质性愚懦不善营运，欠数十万帑银无着落，盐政请将王德法承办常茂号盐产，交通纲商人共出资本选派妥商代为营运。乾隆三十年时，王镗曾旧欠无利之本息等银274,800两，及挪欠永庆号银91,000两，共银365,800两。通纲商人选得商人查奕茂承办常茂号安阳、林县引地，通纲商人具保，该商如有拖欠，通纲代为分赔。查奕茂代办常茂号引课，每年缴帑息银14,069.23余两，分限二十六年归还。79虽说分限十年，实际上常茂号交纳内务府帑息自乾隆三十年到六十年共交了436,146.13两，已经超过他借贷的本银。

2. 两淮盐商借内帑

① 商捐

盐商借内帑的问题已经有许多学者研究，王思治、金成基发表的〈清代前期两淮盐商的盛衰〉提到两淮盐商向皇帝借款，「借帑」实际上成为皇帝所占的干股。80自乾隆十三年起每年商捐银两十万生息，连续五年，共得银84万两。(参见表4），自十八年七月起永远作本生息，余剩银24万两先行解交内库。两淮盐政每年交利息72,000两。81

\begin{BookQuote}
年代银两性质银两(两) 利息共银
\end{BookQuote}

乾隆十三年七月一十四年六月商捐本息 100,000一分五厘 118,000乾隆十四年七月一十五年六月商捐本息 218,000一分五厘 257,000乾隆十五年七月一十六年六月商捐本息 357,000一分五厘 421,000乾隆十六年七月一十七年六月商捐本息 521,000一分五厘 615,000乾隆十七年七月一十八年六月商捐本息 715,000一分五厘 840,000

\begin{BookQuote}
资料来源:<宫中朱批奏折〉，「财政类·盐务」，第0452-0013号，「乾隆十三年六月初七日」
\end{BookQuote}

②盐商借款

由表1可见乾隆前期两淮盐商程可正、程容德、黄源德等商捐达百万两，且被

\vfill
\begin{BookNotes}
\BookNote{79.}{国立故宫博物院编，《宫中档乾隆朝奏折》，第26辑，页416-417。}
\BookNote{80.}{王思治、金成基，〈清代前期两淮盐商的盛衰〉，《中国史研究》，1981年2期，页66-84。}
\BookNote{81.}{〈宫中殊批奏折〉，「财政类·盐务，第0452-0013号，乾隆十三年六月初七日。}
\end{BookNotes}

% 原书第 605 页（PDF 第 618 页）
\BookPage{605}{618}
\noindent 称「徽商向来商力充裕，办运者百数家，有挟资千万者，最少也一、二百万。」82而为江春为商总，「身系两淮盛衰者垂五十年」，为了两淮盐引案「商总江广达等521人联名呈称，商等从前妄用提引余利，罪有应得。」江春于乾隆三十七年起向皇帝借款30万两，年缴利息36,000两，乾隆五十六年两淮盐运使鹿荃文开，众商代办赏借商人江广达运本银300,000两除应交正利银36,000两外，余利银16,000两解交内务府充公。乾隆五十八年两淮盐运使曾燠文开，遵旨赏借总商江广达过继之子江振鸿营运本银5万两，每年按一分起息一年利银6,000两，乾隆六十年江振鸿交利银52,000两。83此外，两淮盐政借商人车脚生息银每年交6,000两。

以上，两淮盐商借内帑的利息大约在10万两左右。周志初在〈两淮盐商「帑利」述略〉统计乾隆年间，清皇室内务府等机构向两淮盐商发放帑本银2,249,756两，乾隆末年两淮盐商每年交纳的帑利大概已在30万两左右。84笔者查阅嘉庆朝《两淮盐法志》借帑之记载，乾隆五十年（1785）乙巳纲起至嘉庆六年（1801）辛酉纲止，该商等历届请借银两惟甲寅（1794）、丙辰（1796）、辛酉等纲领至140余万及160余万不等，此外各纲俱不过120至130万两，并有不及百万两者。85上述盐商借款是否来自内帑并不清楚，〈内务府银库进项月折档〉没有记载这项借款的利息，而且月折档只到乾隆六十年就结束，嘉庆年间盐商如何摊还利息可能必须参考宫中档朱批奏折。

3、内帑的利率

大陆学者李克毅撰写〈清代盐商与帑银〉一文，认为盐商向内务府借钱是因其利息较低，年利率多半是12\%，比民间高利贷年利率30\%低。86不过，他讨论乾隆晚期的帑息超过20\%，他计算方式可能还需讨论。就商人王德宜乾隆五十一年借内务府银300,000两作为帑本，加利银300,000两，本利银计600,000两，分限十五年

\vfill
\begin{BookNotes}
\BookNote{82.}{李澄，《淮鹾备要》，卷7，转引自汪士信，〈乾隆时期徽商在两淮盐业经营中应得、实得利润与流向试析〉，《中国经济史研究》，1989年3期，页96。}
\BookNote{83.}{〈内务府银库进项月折档〉，乾隆四十年八月；五十六年十月；六十年六月；六十年十月。}
\BookNote{84.}{周志初，〈两淮盐商「帑利」述略〉，《盐业史研究》，1991年3期，页58-60。}
\BookNote{85.}{佶山监纂、单渠等总纂，《两淮盐法志》（台北：中央研究院傅斯年图书馆藏嘉庆11年（1806）刊本），卷17，页10。}
\BookNote{86.}{李克毅，〈清代盐商与帑银〉，《中国社会经济史研究》，1989年2期，页19-24。}
\end{BookNotes}

% 原书第 606 页（PDF 第 619 页）
\BookPage{606}{619}
\noindent 带还，每年还本利银40,000两。87李克毅算出利率为22\%，若是这样的利息十五年应该是990,000两，而不是300,000两。故本利分限带还的利率应该为6.66\%。

另一种核算利息的方式称为「利随本减」。长芦盐政西宁文开交乾隆四十二年山东商人内帑本银15万两经内务府大臣议准按月一分起息，分作十五年带还。每年缴本银1万两，利随本减按限解交内库，应解本银1万两，利银18,000两，88利息为12\%。

内务府的利息或许不高，但商人借内帑数量庞大，到乾隆五十七年的时候长芦盐商期借欠的银两相当多。五十七年长芦盐政奏称，商人等节年未交帑项，有应交利银本银5,610,500两，又蒙恩停利归本银2,466,000余两，又摊纲银1,677,248两，又捐饷银500,000两，共银10,018,810两。89此时，长芦盐商的负债额与雍正末年所欠十六万两相比，增加了三十余倍。90

王士铭研究长卢盐政和盐商营运的关系，提到内务府引地租借或帑利借款，对于盐商而言，透过长芦盐政的特殊身份与皇帝沟通，往往获得的帮助是较为快速且显而易见的。91

长芦盐商借内帑到乾隆四十七年的时候已经没办法支付利息，长芦盐政穆腾额奏称，该年因长芦盐商借欠帑本帑利796,900余两，自四十八年起每引缴银0.2两，至五十二年过交银914,390余两，每年约收银200,000两。现在尙有380,000余两，未能还清，请于下届五年再征收0.2两。92盐商欠银却摊在每引中，至此以后，商人更有恃无恐，变本加厉地借内帑。

\vfill
\begin{BookNotes}
\BookNote{87.}{〈内务府银库进项月折档〉，乾隆55年12月。乾隆六十年七月载:乾隆六十年长芦盐政方维甸咨称，将58、59年应交帑利各扣留银40万两借给众商接济运本，仍加利银80万两，限16年完缴。今各商完交，乾隆59年各款利银内除扣留赏借众商银40万两，共银1,600,000两。}
\BookNote{88.}{国立故宫博物院编，《宫中档乾隆朝奏折》，第45辑，页775-776。}
\BookNote{89.}{〈宫中殊批奏折〉，「财政类·盐务」，第0481-010号，乾隆六十年八月十一日。}
\BookNote{90.}{芮和林，〈浅析乾隆时期长芦盐商走向衰落的原因〉，《盐业史研究》，1994年4期，页22-26。}
\BookNote{91.}{王士铭，〈既是官员也是奴才：乾隆朝长芦盐政〉（南投：国立暨南国际大学历史学系硕士论文，2007)，页76-113。}
\BookNote{92.}{国立故宫博物院编，《宫中档乾隆朝奏折》，第68辑，页647-648。}
\end{BookNotes}

% 原书第 607 页（PDF 第 620 页）
\BookPage{607}{620}
盐商积欠内帑被参革，其无着课帑及杂款摊入商捐参课银项下。所谓「参课」是盐商欠项，每引捐银二钱。长芦盐政穆腾额文开，参退各商拖欠无着课帑及杂款银共1,279,938.74两，前经奏准。自乾隆五十四年起每年由征收商捐0.2两参课银两项下，先行拔补帑本次归利银，交五十六年分拨补帑本银71,722.96两。93

乾隆五十九年长芦盐政征瑞奏称，长芦向有参商引窝家产不敷变抵欠帑，于乾隆四十七年初任长芦查出积年无着银790,000两，奏请自四十八年为始，通纲每引捐银0.2两，经内务府议覆准行。旋因逐年复有参商无着银两，又经奏准商捐0.2两永远征收。又有各参商案续入者，通共未完银1,279,900余两。自五十三至五十七止共收0.2两商捐855,299两，除前盐政将五十三年所收先行拨补参课银133,900余两报明内务府外，实存银721,383两。此内拨出经费必须转发别商生息银123,211两，又抵参商牛绳祖等带征银132,583两，又拨出帑本解内务府银182,704两，三款俱照案拨补。又发过剥船津贴饭银135,090两，解过香灯现年利银47,058两。并续参商章惠生、卫国泰等无着帑课先行拨抵银100,734两，共银721,283两，俱逐年造报内务府查核外，所有五年商捐0.2两一项全数拨抵净未完银841,438两。参商章惠生、卫国泰等无着帑课共466,260余两，除先拨抵外尚未拨银365,534两。又查出参商王致玉等到限缓课银23,175两，共无着银388,709两。请自六十年起每引随课0.05两共0.25两，核计每年可以多收银五万余两。以此抵剥船岁须津贴银27,000两。94

\BookSubheading{四、其他收入}

长芦商人交节省银始自乾隆十二年。雍正朝李天馥充商以后应交引价并未完，铜斤节省尚未届限，迨至雍正九年因欠课革退，招商范宗文接办引窝任交帑项。李天馥认办张霖盐窝止认变抵银165,800余两，当时李天馥原因已身欠有铜斤节省银111,300两，故呈请俟引价铜斤完日，情愿每年再交二成节省银3,620两。范宗文虽将认纳帑银俱各按限完结，故交长芦禹临等八处引窝二成节省银，每年

\vfill
\begin{BookNotes}
\BookNote{93.}{〈宫中朱批奏折〉，「财政类·盐务」，第0477-047号，乾隆五十七年闰四月二十三日。}
\BookNote{94.}{〈宫中朱批奏折〉，「财政类·盐务」，第0480-028号，乾隆五十九年十月十五日。}
\end{BookNotes}

% 原书第 608 页（PDF 第 621 页）
\BookPage{608}{621}
\noindent 交3,620两。95

据长芦盐政丽柱奏称，长芦各商普加盐斤，内有河南祥符等州县引地系办理官引商人王瑶等承办，亦经一体加给盐斤，在于免交钱粮之内，在民商所获余利无庸置议。至王瑶现办引地从前并未办公效力，且以抵归内务府官项，而加斤之恩在王瑶为便滥邀，王瑶禀称承办引地系属官项，每包节省额课银0.64两，按库平足色共应交5,451两，请照数归公。96王瑶承办官引地，于每包银内课节省银归公，此项银两系交内务府。王至德每年交送盐引节省银为4,995.43两。若商人被革退，代办的商号必须承担这项节省银，如义和泰接手经营王至德盐引地，每年仍须缴交节省银4,995.43两。王德宜代退商范清济办引地亦交节省银。

其次是引地的地租。满城县官引地除了交正额税课外，每引按5两成本合计起租，原本是郭有恒办理，当他被革退后由张益昌继任，每年交银3,022.5两。

乾隆十四年两淮盐政吉庆谨奏:淮北商人程容德寄顿张广泗银两一案，程容德情愿捐银30,000两赎罪自新。97乾隆五十七年两淮盐运使柴桢侵挪库银22万两，盐商加倍罚银44万两。此案系因盐商上纳钱粮，或用商伙或用家人赴库兑交，本商并不侵到。其小商并有交与银店代纳者，此次该运使以需用甚急，又不便将入库之项再行提出，有家柏顺住居署旁代该运使设计在外揽收，向交纳钱粮之商伙家人，及银店人等说明，将应纳钱粮交伊收受，即塡入库簿作为已收。各商伙家人及银店人等利其便于缴纳，且银色稍有不足可免倾销火耗，其中不无些微沾润，遂皆向伊交收。98有关柴桢案件的档案相当多，等将来有机会再仔细讨论。

五结论

刘翠溶教授认为内务府和户部的关系，一者管皇室财政；一者管国家财政，然而两者权限并未绝对的划分，经费也互有往来，其间关系逐渐由权宜演变成定

\vfill
\begin{BookNotes}
\BookNote{95.}{〈宫中殊批奏折〉，「财政类·盐务」，第0451-027号，乾隆十二年四月二十七日。}
\BookNote{96.}{〈宫中朱批奏折〉，「财政类·盐务」，第0453-028号，乾隆十四年十一月二十日。}
\BookNote{97.}{〈宫中殊批奏折〉，「财政类·盐务」，第0453-024号，乾隆十四年十一月十五日。}
\BookNote{98.}{〈宫中朱批奏折〉，「财政类·盐务」，第0478-029号，乾隆五十七年十二月十日。}
\end{BookNotes}

% 原书第 609 页（PDF 第 622 页）
\BookPage{609}{622}
\noindent 例。99MadeleineZelin（曾小萍）则认为满族王朝倾力进行财政改革，以强化君主专制并加强中央集权政府的各项制度建设。他们将皇室内府与国库分开，反映了内廷和外朝有着明晰的划分。100由本文讨论可知，户部和内务府财政并非截然划分，尤其乾隆朝的盐税方面，很明显地内务府透过盐政的运作，不断地侵蚀户部的财政。其次，陈国栋认为内务府包衣充任税差间接干预国家财政，对于拓展内务府财政的贡献实属有限，反而是多充监视地方耳目，专制统治比较实际。101然而，乾隆皇帝说「各省盐政、织造、关差皆系内府世仆」，盐政和关差都是替皇帝荷包攒添了大量财富。

据乾隆朝《大清会典》记载，乾隆十八年（1753）的税赋，盐税为5,560,540两。至嘉庆十七年(1812)盐税为5,797,645两。102历经清朝盛世，国家的重要税收支一的盐税似乎没有增加。反观皇室财政，由乾隆初的数十万两到乾隆晚期达一、二百万两。从本文的讨论可知，乾隆初期原本交给户部的盐政衙门经费，改解内务府；衙门的运库银两借给商人，利息却归皇帝所有；盐商欠税课，查抄家产也归皇帝所有。因此，户部与皇室财政的消长可以看出盐商对于皇室财政实质的贡献。

许多学者认为盐商对皇帝的捐输和报效为「巨额开支」，本文却从不同角度观察。第一、两淮盐政衙门的经费给予内务府的收入共18,811,258.95两，大约占盐政官员所得95,482,800两之20\%。盐商除了正课之外，私费不赀，「遇一事即有一事之陋规，经一处即有一处之科派」，各项陋规或科派，皇帝也获得不少好处。长芦

\vfill
\begin{BookNotes}
\BookNote{99.}{刘翠溶，《顺治康熙年间的财政平衡问题》（台北：嘉新水泥公司文化基金会丛书，1969），页110}
\BookNote{100.}{曾小萍着，董建中译，《州县官的银两：十八世纪中国合理化财政改革》（北京：中国人民大学出版社，2005)，页4。}
\BookNote{101.}{陈国栋，〈清代中叶以后重要税差专由内务府包衣担任的几点解释〉，收入许倬云、毛汉光、刘翠溶主编《第二届中国社会经济史研讨会论文集》（台北：汉学研究资料及服务中心，1983)，页173-204。}
\BookNote{102.}{允祹等纂修，乾隆《大清会典》（光绪34年上海商务印书馆印本，中研院近史所图书馆藏线装书)，卷10，页2；卷15，页3；卷16，页1。《汇核嘉庆十七年各直省钱粮出入清单》（1930年刊行），转引自香昌纪撰，赵中男译，〈论清前期嘉庆年间的国家财政与关税收入〉，《社会科学辑刊》，1993年3期，页88-94。}
\end{BookNotes}

% 原书第 610 页（PDF 第 623 页）
\BookPage{610}{623}
\noindent 盐政衙门的经费可能是借给盐商，故其衙门经费缴交内务府的数量不及两淮盐政。

第二、学者们统计乾隆朝盐商报效军需、河工、皇帝南巡达三千万两，实际上，从〈月折档〉所见盐商对皇帝的报效不过15次，这些负担不及提引余利或盐商长期借帑的利息。两淮盐商报效大约在乾隆十一年实施提引制度后次数较多，而且也替乾隆皇帝六次南巡，耗资巨万建造行宫园林，至乾隆三十三年两淮盐引案发，盐商赔补六百多万两，又缴交提引余利四百多万两，造成商人资金短绌，连总商江春都跟皇帝借内帑。

第三、长芦盐商和两淮盐商借内帑的比例，以长芦盐商借帑较多，到乾隆五十七年时长芦盐商期借欠的银两相当多，盐商欠本利银达千万两以上。长芦盐商借欠帑本帑利，除了按月分期付款外，自乾隆四十八年起到五十七年摊入盐引中，每引缴银0.2两，每年约收银200,000两。将盐商借钱摊入盐引税课中，是不合理的方式。可是，皇帝谕旨有谁能抗衡？两淮盐商摊还的提引余利还不是同样的遭遇。

皇帝将内帑借给盐商利润丰厚，不过，到咸丰年间却成为呆帐，根据光绪年间内务府银库进款档记载，两淮盐商借款利息364,200两自咸丰元年（1851）至今未交，长卢盐商借款利息104,800两，近年约交六、七万两。103咸丰元年内务府丧失帑利收入，以致于后来由户部由盐课、税关中添拨经费，形成户部与内务府财政紧张关系。

汪崇赏在清嘉道时期两淮官盐的壅滞〉一文中，认为皇帝对「商人的感情倾注和体恤，恐怕也是很特殊的。尤其是徽州大盐商江春，堪称『以布衣上交天子』的典范」。104盐商众多的支出而走向没落，盐商「上交天子」恐怕是哑巴吃黄莲有苦说不出。

\vfill
\begin{BookNotes}
\BookNote{103.}{〈内务府银库进款档〉（北京：中国第一历史档案馆藏），档案编号1678，光绪元年二月。}
\BookNote{104.}{汪崇，〈清嘉道时期两淮官盐的壅滞〉，《盐业史研究》，2002年第4期，页3-9。}
\end{BookNotes}

% 原书第 611 页（PDF 第 624 页）
\BookPage{611}{624}
\BookArt{p624-01.jpg}{0.990}{0.452}


% 原图说　图6　乾隆时期盐政总收入

% 原书第 612 页（PDF 第 625 页）
\BookPage{612}{625}
